% Diubah pada 4/2/97.
Teori peluang berawal di Prancis pada abad ketujuh belas ketika dua
matematikawan besar Prancis, Blaise Pascal dan Pierre de Fermat, berkirim surat
tentang dua masalah dalam permainan untung-untungan. Masalah seperti yang
diselesaikan Pascal dan Fermat terus memengaruhi para peneliti awal seperti
Huygens, Bernoulli, dan DeMoivre dalam membangun teori matematis tentang
peluang. Kini, teori peluang merupakan cabang matematika yang mapan dan
diterapkan dalam setiap bidang kegiatan keilmuan, dari musik hingga fisika,
serta dalam kehidupan sehari-hari, dari prakiraan cuaca hingga perkiraan risiko
pengobatan medis baru.

Buku ini dirancang untuk mata kuliah pengantar peluang bagi mahasiswa tahun
kedua, ketiga, dan keempat dalam bidang matematika, ilmu fisika dan sosial,
rekayasa, serta ilmu komputer. Buku ini menyajikan pembahasan menyeluruh
tentang gagasan dan teknik peluang yang diperlukan untuk memperoleh pemahaman
yang kokoh mengenai bidang ini. Buku ini dapat digunakan dalam mata kuliah
dengan beragam durasi, tingkat, dan fokus.

Untuk mata kuliah standar selama satu semester yang mencakup peluang diskret
dan kontinu, mahasiswa sebaiknya telah menempuh dua semester kalkulus,
termasuk pengantar integral lipat. Untuk membahas Bab~\ref{chp 11}, yang
memuat materi tentang rantai Markov, diperlukan pula pengetahuan tentang teori
matriks.

Buku ini juga dapat digunakan dalam mata kuliah peluang diskret. Materinya
disusun sedemikian rupa sehingga pembahasan peluang diskret dan kontinu
disajikan secara terpisah, tetapi sejajar. Susunan ini menghindarkan penyajian
peluang yang terlalu ketat atau formal dan memberikan manfaat pedagogis yang
kuat: pembahasan diskret kadang dapat memotivasi pembahasan peluang kontinu
yang lebih abstrak. Untuk mata kuliah peluang diskret, mahasiswa sebaiknya
telah menempuh satu semester kalkulus.

Hanya sedikit latar belakang komputasi yang diasumsikan atau diperlukan untuk
memanfaatkan sepenuhnya materi dan contoh komputasi dalam buku ini. Semua
program yang digunakan dalam buku ini telah ditulis dalam masing-masing bahasa
TrueBASIC, Maple, dan Mathematica.

Buku ini didistribusikan melalui Web sebagai bagian dari The CHANCE Project, yang
bertujuan menyediakan materi bagi mata kuliah tingkat awal dalam peluang dan
statistika. Program komputer, penyelesaian latihan bernomor ganjil, dan daftar
ralat terkini juga tersedia di situs tersebut. Pengajar dapat memperoleh
seluruh penyelesaian dengan menulis kepada salah satu penulis di
jlsnell@dartmouth.edu atau cgrinst1@swarthmore.edu.

\bigskip\centerline{\bf FITUR}\medskip

{\it Tingkat ketelitian dan penekanan:} Peluang merupakan bidang matematika
yang sangat intuitif dan mudah diterapkan. Kami berusaha agar keindahannya
tidak tertutupi oleh terlalu banyak matematika formal. Sebaliknya, kami
berusaha mengembangkan gagasan-gagasan utama dengan alur yang tidak tergesa,
memberikan beragam penerapan peluang yang menarik, dan menunjukkan sejumlah
contoh yang berlawanan dengan intuisi sehingga peluang menjadi bidang yang
begitu hidup.

{\it Latihan:} Buku ini memuat lebih dari 600 latihan yang memberikan banyak
kesempatan untuk melatih keterampilan dan membangun pemahaman yang kokoh
tentang gagasan-gagasannya. Kumpulan latihan mencakup latihan rutin yang dapat
dikerjakan dengan atau tanpa komputer serta latihan yang lebih teoretis untuk
memperdalam pemahaman konsep-konsep dasar. Latihan yang lebih sulit ditandai
dengan tanda bintang. Buku petunjuk penyelesaian untuk seluruh latihan tersedia
bagi pengajar.

{\it Catatan historis:} Peluang tingkat pengantar merupakan bidang yang
gagasan-gagasan dasarnya masih berkaitan erat dengan gagasan para perintisnya.
Karena itu, buku ini memuat banyak catatan historis, terutama yang berkaitan
dengan perkembangan peluang diskret.

{\it Pemanfaatan pedagogis program komputer:} Teori peluang membuat prediksi
tentang percobaan yang hasilnya bergantung pada keacakan. Karena itu, komputer
sangat cocok digunakan sebagai alat matematis untuk menyimulasikan dan
menganalisis percobaan acak.

Dalam buku ini, komputer digunakan dengan beberapa cara. Pertama, komputer
menyediakan laboratorium untuk menyimulasikan percobaan acak sehingga
mahasiswa dapat mengenali keragaman percobaan tersebut. Pemanfaatan komputer
dalam peluang ini telah digambarkan dengan sangat baik oleh William Feller
dalam edisi kedua buku terkenalnya, {\it An Introduction to Probability Theory
and Its Applications} (New York: Wiley, 1950). Dalam prakata, Feller menjelaskan
bahwa fluktuasi dalam pelemparan koin dapat sedemikian bertentangan dengan
intuisi umum hingga kolega yang berpengalaman pun meragukan prediksi teorinya;
karena itu ia menyertakan catatan percobaan simulasi. Ini merupakan ringkasan
editorial; ungkapan sumber tidak diterjemahkan atau direproduksi.

Selain menyediakan laboratorium bagi mahasiswa, komputer merupakan sarana yang
kuat untuk memahami hasil-hasil dasar teori peluang. Sebagai contoh, gambaran
grafis mengenai pendekatan distribusi binomial terbakukan terhadap kurva normal
memberikan penjelasan yang lebih meyakinkan tentang Teorema Limit Pusat
daripada banyak bukti formal untuk hasil mendasar tersebut.

Terakhir, komputer memungkinkan mahasiswa menyelesaikan masalah yang tidak
memiliki rumus bentuk tertutup, misalnya waktu tunggu dalam antrean. Bahkan,
kehadiran komputer mengubah cara kita memandang banyak masalah dalam peluang.
Sebagai contoh, kemampuan menghitung peluang binomial secara tepat untuk
percobaan dengan hingga 1000 uji coba mengubah cara kita memandang pendekatan
normal dan Poisson.

\bigskip\centerline{\bf UCAPAN TERIMA KASIH (ACKNOWLEDGEMENTS)}\medskip

Siapa pun yang menulis buku peluang pada masa kini berutang budi besar kepada
William Feller, yang mengajari kita semua cara menghidupkan peluang sebagai
sebuah bidang kajian. Jika Anda menemukan contoh, penerapan, atau latihan yang
sangat Anda sukai, kemungkinan besar asalnya adalah buku klasik Feller,
{\sl An Introduction to Probability Theory and Its Applications.}

Kami berutang budi kepada banyak orang atas bantuan mereka dalam pekerjaan
ini. Pendekatan terhadap rantai Markov yang disajikan dalam buku ini
dikembangkan oleh John Kemeny dan penulis kedua. Reese Prosser menjadi rekan
penulis yang tidak disebutkan untuk materi peluang kontinu dalam versi
terdahulu buku ini. Mark Kernighan menyumbangkan 40 halaman komentar mengenai
edisi sebelumnya. Banyak di antara komentar tersebut sangat menggugah
pemikiran; selain itu, komentar-komentar itu memberikan sudut pandang mahasiswa
terhadap buku ini. Sebagian besar perubahan utama dalam versi ini berawal dari
catatan tersebut.

Fuxing Hou dan Lee Nave memberikan banyak bantuan dalam tata letak dan gambar.
John Finn memberikan saran pedagogis yang berharga tentang teks dan program
komputer. Karl Knaub dan Jessica Sklar bertanggung jawab atas implementasi
program komputer dalam Mathematica dan Maple. Jessica dan Gang Wang membantu
menyusun penyelesaian.

Terakhir, kami berterima kasih kepada American Mathematical Society, khususnya
Sergei Gelfand dan John Ewing, atas perhatian mereka terhadap buku ini, bantuan
mereka dalam produksinya, dan kesediaan mereka menjadikan karya ini dapat
didistribusikan ulang secara bebas.
